
% \vspace{-0.1in}

\section{Related Work}
\label{related_work}

%\vspace{-0.025in}

\textbf{Accelerator for Transformer Model}
Hardware accelerators that support transformer-based NLP models have been recently proposed. A3 \cite{ham2020a3}, GOBO \cite{zadeh2020gobo}, SpAtten \cite{wang2021spatten}, EdgeBERT \cite{tambe2021edgebert}, and ELSA \cite{ham2021elsa} discuss designs that only speed up the attention mechanism in the transformer using various pruning, quantization, or both, without taking into consideration the end-to-end process. For instance, SpAtten uses pruning to accelerate the attention mechanism and the layer normalization process but does not consider token embedding, residual, and LM head. Our work targets the language services at the datacenter, especially focusing on text generation workloads, so there is less emphasis on aggressively speeding up a specific operation in the transformer but more emphasis on the model-parallel dataflow that can address both memory and compute-intensive problems in its end-to-end inference.

\textbf{Hardware Architecture for Datacenters} Many domain-specific acceleration architectures have been proposed for machine learning but few at the scale of datacenters. Microsoft's Brainwave \cite{fowers2018configurable} is a relevant work that implements an FPGA-based neural processing unit for datacenters with high-performance processing units. However, Brainwave does not utilize high-bandwidth off-chip memory, so it cannot effectively run memory-intensive workloads, such as GPT inference. Google \cite{jouppi2017datacenter, jouppi2020domain, jouppi2021ten} proposes the systolic-based TPU architecture for various DNN applications, and Facebook \cite{anderson2021first} designs accelerators specialized for recommendation systems. A few of these previous architectures for datacenters \cite{fowers2018configurable, jouppi2021ten} support the transformer, but none have design optimizations to drastically improve the performance for text generation workloads specifically.

%None of these previous architectures for datacenters support the transformer because the transformer has processes like attention that have different operational characteristics than other DNNs.
